\label{sec:spectrum}

Papers M.1 through M.7 each implement one community character signal.
Table~\ref{tab:signal-spectrum} provides an overview of all seven signals: their
data source, the similarity metric used, geographic availability, and the legal
precedent most directly supporting their use as a community of interest proxy.

\begin{table}[ht]
\centering
\caption{Spectrum of community character signals across the M-track.
  All sources are free and publicly available from federal agencies.
  Coverage refers to fraction of US states with usable data.}
\label{tab:signal-spectrum}
\small
\begin{tabularx}{\textwidth}{@{}lXlll@{}}
\toprule
\textbf{Paper} & \textbf{Signal} & \textbf{Data Source} & \textbf{Similarity} & \textbf{Coverage} \\
\midrule
M.1 & Economic character & LODES WAC & Cosine & 100\% \\
M.2 & Land use mix & NLCD & Cosine & 100\% \\
M.3 & Housing character & ACS & Cosine & 100\% \\
M.4 & Commuting shed & LODES OD & Jaccard & 100\% \\
M.5 & Topographic zone & USGS 3DEP & Valley/ridge & 100\% \\
M.6 & Zone co-membership & TIGER/EIA 861 & Count ratio & 100\%$^a$ \\
M.7 & Transit accessibility & GTFS & Isochrone overlap & $\sim$70\% \\
\bottomrule
\end{tabularx}
\smallskip
\begin{minipage}{\textwidth}
  \small $^a$ School district + county subdivision + electric utility provide
  100\% national coverage. Fire and water districts cover approximately 60\%
  and 40\% of states, respectively.
\end{minipage}
\end{table}

\subsection{M.1 --- Economic Character (LODES WAC)}

The Longitudinal Employer-Household Dynamics (LEHD) Origin-Destination Employment
Statistics (LODES) Workplace Area Characteristics (WAC) file provides, for each
census block, counts of jobs in 20 industry categories, 3 earnings bands,
and 8 worker age groups~\citep{lodes2022}.
Aggregated to the tract level, this produces an occupational character vector
reflecting where a tract fits in the regional economy: a manufacturing district,
a financial services node, a retail corridor, a hospital campus.
Two tracts with similar industry mixes are, economically, members of the same
community of interest.
Cosine similarity of the industry-share vector is the natural measure: it rewards
proportional similarity of economic composition regardless of the absolute scale
of employment.

Economic communities of interest are explicitly recognized by Colorado's constitution
(``trade area'' factors) and have been accepted as a valid community type in Arizona
IRC proceedings.
LODES data is annual, free, and available from the Census Bureau's LEHD program.

\subsection{M.2 --- Land Use Mix (NLCD)}

The National Land Cover Database (NLCD), produced by the Multi-Resolution Land
Characteristics Consortium, classifies every 30-meter pixel in the United States
into one of 20 land cover categories~\citep{nlcd2019}.
Aggregated to the tract level, the NLCD produces a land use composition vector
(fraction of developed land, fraction of forest, fraction of agricultural land,
fraction of wetland, etc.) that captures the physical character of the tract.
Two urban tracts are similar; two agricultural tracts are similar; an urban tract
and an agricultural tract are dissimilar.
The NLCD is updated approximately every five years and is freely available.

Physical geography as a community character signal has ancient legal roots.
Natural features---rivers, mountain ranges, valley floors---have served as district
boundaries since the colonial era, and state constitutions in many western states
explicitly mention ``natural boundaries'' as a criterion.
The NLCD operationalizes physical character systematically, across the entire landscape,
without requiring the map-drawer to identify specific named features.

\subsection{M.3 --- Housing Character (ACS)}

The American Community Survey (ACS) five-year estimates provide, at the tract level,
counts of housing units by structure type (single-family, multi-family, mobile home),
tenure (owner-occupied, renter-occupied), vacancy status, and year built~\citep{acs2020}.
The housing character vector captures the residential identity of a tract: an
established owner-occupied suburb of single-family homes built in the 1960s is a
different community from a high-density rental market with recently constructed
apartment buildings.
Cosine similarity of the housing character vector captures community similarity
in the sense of ``neighbors who share a stake in the same type of housing market.''

Housing character is particularly relevant to the communities of interest doctrine
because housing policy---zoning, rent control, historic preservation, homeowners
association governance---is a primary vehicle through which neighbors exercise
collective civic interests.
Homeowners in the same zoning district share an interest in land-use decisions;
renters in the same rental market share an interest in rent stabilization policy.

\subsection{M.4 --- Commuting Shed (LODES OD)}

The LODES Origin-Destination (OD) file records, for each pair of census blocks
with commuting flows, the count of workers who live in one block and work in
another~\citep{lodes2022}.
Aggregated to the tract level, this produces a destination distribution for each
residential tract: where do its residents go to work?
Two tracts that send large fractions of their workers to the same employment centers
are parts of the same functional economic area---the commuting shed or labor market area.
Jaccard similarity of the destination set (tracts that receive at least one commuter
from each origin) captures this relationship.

The commuting shed is a well-established concept in regional science and is recognized
as a proxy for economic community of interest by California's commission guidance
and Colorado's ``trade area'' criterion.
Functional labor markets are a natural unit of democratic representation: residents
who commute together to the same city have shared interests in transportation
infrastructure, highway policy, and regional economic development.

\subsection{M.5 --- Topographic Zone (USGS 3DEP)}

The USGS 3D Elevation Program (3DEP) provides one-meter to one-third-arc-second
digital elevation models covering the contiguous United States~\citep{usgs3dep}.
From elevation data, the topographic zone of a tract can be classified as valley
floor, upland plateau, ridge, or slope zone using standard terrain analysis methods
(slope curvature and hydrological drainage basin algorithms).
Two tracts in the same valley share a drainage basin, a flood risk profile, an
agricultural character, and often a historical settlement pattern; two ridge tracts
share a different set of interests.

Topographic community signals are particularly relevant in western states, where
valley communities in the Sierra Nevada, Cascades, Rockies, and Wasatch Range are
economically and culturally distinct from adjacent upland areas.
Natural boundaries have been cited approvingly in redistricting cases as evidence
that a district boundary is not arbitrary~\citep{ncsl2021criteria}.

\subsection{M.6 --- Administrative Zone Co-membership (TIGER/EIA 861)}

Zone co-membership is the most legally direct community character signal.
When two tracts share a school district, fire district, water district, electric
utility service territory, and county subdivision, they are members of the same
administrative community in every governmentally recognized sense.
The zone score
\begin{equation}
  \mathrm{zone\_score}(u,v) \;=\; \frac{|\{z \in Z : \mathrm{zone}_z(u) = \mathrm{zone}_z(v)\}|}{|Z|}
  \label{eq:zone-score}
\end{equation}
counts how many of the available zone types are shared, normalized to $[0,1]$.
The zone co-membership formula replaces cosine similarity in equation~(\ref{eq:weight-modifier})
and satisfies all required properties: symmetry, $[0,1]$ range, and graceful degradation
to $w_{\mathrm{base}}$ when no zone data is available.
Data for school districts and county subdivisions comes from TIGER/Line shapefiles
(100\% national coverage); electric utility service territories come from EIA
Form 861 (100\% coverage); fire and water districts from TIGER/Line Special District
files (\textasciitilde 60\% and \textasciitilde 40\% coverage, respectively).
M.6 provides full documentation of the zone co-membership algorithm and its legal grounding.

\subsection{M.7 --- Transit Accessibility (GTFS)}

General Transit Feed Specification (GTFS) data, published by transit agencies and
aggregated by Transitland and the National Transit Database~\citep{ntd2022},
encodes transit routes, stops, and schedules for approximately 70\% of the US
by population.
Transit accessibility can be measured as the set of destinations reachable within
a fixed travel time (e.g., 45 minutes) from a tract centroid; two tracts with
overlapping reachable-destination sets are members of the same transit-accessible
community.
Isochrone overlap---the fraction of each tract's reachable set that is also reachable
from the other tract---provides a natural similarity measure.

Transit co-membership is a recognized community of interest in California's commission
guidance and in urban planning literature on transit-oriented development.
In dense urban areas, transit accessibility is a primary determinant of where
residents can work, shop, and access services, making it a more meaningful community
bond than road network proximity alone.
M.7 is limited to the approximately 70\% of states and counties with published GTFS
feeds; the algorithm degrades gracefully to the zone co-membership or geographic
baseline for jurisdictions without transit data.

\subsection{M.8 --- Composite Community Index (preview)}

Paper M.8, not yet published, will define a composite community character index
that combines M.1--M.7 signals with expert-tunable weights.
The composite is a weighted average of the individual similarity scores,
providing a single number that can be used as the edge weight modifier when all
signals are active.
M.8 will also provide an expert witness template: a standardized declaration format
that maps the composite score to specific communities of interest recognized in
prior redistricting proceedings.
